% Book 4, Year 2, chapter 8 -- Traits and generics.
% Provenance (never printed): sources/book4/BRIEFS.md, chapter 8.
\chapter{Traits and Generics}\label{ch:b2:rust-traits-and-generics}

\code{println!("{}", x)} prints integers, floats, strings, and, after five lines of code, a type of our own that the
authors of \code{println!} never saw.

\omcode{code/bachelor-2/08-rust-traits-and-generics/listings/display.rs}{4}{21}{A book that prints
itself}\label{lst:b2:rust-traits-and-generics:display}

\omoutput{code/bachelor-2/08-rust-traits-and-generics/outputs/display.txt}{One method written, three ways to print.
The width 30 had no effect.}\label{out:b2:rust-traits-and-generics:display}

\code{println!}, \code{format!} and \code{to_string} ask one thing of their argument: that its type implements the
\omterm{def:b2:rust-traits-and-generics:trait}{trait} \code{Display}, a one-method interface. This chapter is about such interfaces, the functions written once for
every type that has them, and the two ways, by copies or by tables, in which the compiler turns them into machine
code.

\begin{recall}
\Cref{ch:b2:cpp-polymorphism}: a \omterm{def:b2:cpp-polymorphism:virtual}{virtual function} is called through the object's pointer to its class's table, at a
cost of one indirect call. \Cref{ch:b2:cpp-templates}: a \omterm{def:b2:cpp-templates:template}{template} is instantiated once per argument type, and a
concept states what the argument must provide. \Cref{ch:b2:rust-enums-and-errors}: a derive attribute writes
\code{Debug}, \code{PartialEq} and orders for a type.
\end{recall}

\section{Traits}\label{sec:b2:rust-traits-and-generics:traits}

\begin{definition}[Trait, trait implementation, default method]\label{def:b2:rust-traits-and-generics:trait}
A \emph{trait}\index{trait} is a named set of method signatures, and possibly of \omterm{def:b2:rust-traits-and-generics:associated}{associated types} and constants, that
a type may provide. A \emph{trait implementation}\index{trait implementation}, \code{impl Trait for Type { ... }},
gives the bodies for one type. A \emph{default method}\index{default method} is a method whose body is written in
the trait itself, in terms of the other methods; an implementation may keep it or replace it.
\end{definition}

\omcodes{code/bachelor-2/08-rust-traits-and-generics/listings/bounds.rs}{4-29}{A \omterm{def:b2:rust-traits-and-generics:trait}{trait} for what the library
lends}\label{lst:b2:rust-traits-and-generics:lendable}

\code{Book} keeps the default \code{describe}; \code{Film} replaces it. Unlike a C++ \omterm{def:b2:cpp-polymorphism:inheritance}{base class}, a \omterm{def:b2:rust-traits-and-generics:trait}{trait} has no fields
and is not a type of values; and an implementation may be written after both the type and the \omterm{def:b2:rust-traits-and-generics:trait}{trait}, in another
file, as \code{Display} was for \code{Book}. One restriction keeps this sound.

\begin{definition}[Coherence, orphan rule]\label{def:b2:rust-traits-and-generics:coherence}
\emph{Coherence}\index{coherence (orphan rule)} requires that a program have at most one implementation of a given
\omterm{def:b2:rust-traits-and-generics:trait}{trait} for a given type. Rust enforces it by the \emph{orphan rule}: a crate may write \code{impl Trait for Type} only
if the \omterm{def:b2:rust-traits-and-generics:trait}{trait} or the type is defined in that crate.
\end{definition}

So our crate may implement \code{Display} for \code{Book} (our type) or \code{Lendable} for \code{String} (our \omterm{def:b2:rust-traits-and-generics:trait}{trait}),
but not \code{Display} for \code{Vec<u32>}: two crates could each write one, and the program would not know which to
call. The usual way round is a \emph{\omterm{def:b2:rust-traits-and-generics:newtype}{newtype}} (\cref{sec:b2:rust-traits-and-generics:zero}).

\section{Generic functions and trait bounds}\label{sec:b2:rust-traits-and-generics:generics}

\begin{definition}[Generic function, trait bound]\label{def:b2:rust-traits-and-generics:bound}
A \emph{generic function}\index{generic function} has type parameters, \code{fn f<T>(...)}, and is checked once,
for every type its parameters may take. A \emph{trait bound}\index{trait bound} \code{T: Trait} restricts \code{T} to
types that implement \code{Trait}, and is the only source of what the body may do with a \code{T}. Bounds are
written in the angle brackets, in a \code{where} clause, or as \code{impl Trait} in a parameter's type.
\end{definition}

\omcode{code/bachelor-2/08-rust-traits-and-generics/listings/bounds.rs}{31}{49}{Three ways to write a
bound}\label{lst:b2:rust-traits-and-generics:bounds}

\omoutput{code/bachelor-2/08-rust-traits-and-generics/outputs/bounds.txt}{\code{label} accepts a \code{char} and an
integer as the shelf: both are \code{Display}.}\label{out:b2:rust-traits-and-generics:bounds}

A body that uses an operation its bounds do not promise is refused, before any call is seen:

\omcode{code/bachelor-2/08-rust-traits-and-generics/listings/missing_bound.rs}{3}{10}{A comparison with no
bound}\label{lst:b2:rust-traits-and-generics:missing}

\omoutput[2-13]{code/bachelor-2/08-rust-traits-and-generics/outputs/missing_bound.txt}{The compiler names the missing
\omterm{def:b2:rust-traits-and-generics:trait}{trait}.}\label{out:b2:rust-traits-and-generics:missing}

This is the difference with C++ \omterm{def:b2:cpp-templates:template}{templates}. A \omterm{def:b2:cpp-templates:template}{template}'s body is checked only when it is instantiated, against the
actual type, and an error inside it surfaces at the caller, often deep in a library (\cref{ch:b2:cpp-templates});
a concept may state requirements, but the body is not checked against them. A Rust generic body is checked once,
against its bounds: if it compiles, it compiles for every type that satisfies them, and a caller's error is only
ever \emph{this type does not implement that \omterm{def:b2:rust-traits-and-generics:trait}{trait}}.

\begin{definition}[Monomorphisation]\label{def:b2:rust-traits-and-generics:mono}
\emph{Monomorphisation}\index{monomorphisation} is the compilation of a \omterm{def:b2:rust-traits-and-generics:bound}{generic function} into one machine-code copy
for each combination of type arguments with which the program uses it, each copy compiled and optimised as if it had
been written for those types.
\end{definition}

\omcode{code/bachelor-2/08-rust-traits-and-generics/listings/mono.rs}{2}{18}{One generic, two
uses}\label{lst:b2:rust-traits-and-generics:mono}

\omoutput{code/bachelor-2/08-rust-traits-and-generics/outputs/run_mono.txt}{The symbols of the object file:
\code{largest} exists twice.}\label{out:b2:rust-traits-and-generics:mono}

\begin{proposition}[One copy per type]\label{prop:b2:rust-traits-and-generics:mono}
Let a program call a \omterm{def:b2:rust-traits-and-generics:bound}{generic function} \code{f<T>}, directly or through other \omterm{def:b2:rust-traits-and-generics:bound}{generic functions}, with the set
$\mathcal{T}$ of type arguments, once all generic calls are expanded. Then the compiled program contains at most
$|\mathcal{T}|$ copies of \code{f}, one per element of $\mathcal{T}$, whatever the number of call sites.
\end{proposition}

\begin{proof}
The compiler generates a copy of \code{f} for a type argument when it first meets a call with it, and records it;
every later call with the same argument, from any site, uses the recorded copy. A copy is made only for an argument
that some call uses, so the copies are in bijection with a subset of $\mathcal{T}$: at most $|\mathcal{T}|$. (Fewer
when the optimiser inlines a copy everywhere and drops it, or merges two identical copies, as the weekend problem
will show.)
\end{proof}

The copies are fast, since each is specialised: \code{largest::<u32>} compares with an integer instruction, the
\code{f64} one with a floating-point one, with no indirection. They also cost space and compile time, in proportion
to $|\mathcal{T}|$.

\section{Trait objects}\label{sec:b2:rust-traits-and-generics:dyn}

A vector holds values of one type. To hold rectangles and circles together, the program needs one type for
\emph{any value that implements \code{Shape}}.

\begin{definition}[Trait object, fat pointer]\label{def:b2:rust-traits-and-generics:trait-object}
A \emph{trait object}\index{trait object} \code{dyn Trait} is a value of some type implementing \code{Trait}, whose
type is known only at run time. It is used behind a pointer, \code{&dyn Trait} or \code{Box<dyn Trait>}, which is a
\emph{fat pointer}\index{fat pointer}: the address of the value together with the address of a table, the
\emph{\omterm{def:b2:cpp-polymorphism:vtable}{vtable}}, holding the value type's size, alignment, drop function and the addresses of its implementations of
the \omterm{def:b2:rust-traits-and-generics:trait}{trait}'s methods.
\end{definition}

\omcode{code/bachelor-2/08-rust-traits-and-generics/listings/shapes_dyn.rs}{20}{34}{Two types in one
vector}\label{lst:b2:rust-traits-and-generics:shapes}

\omoutput{code/bachelor-2/08-rust-traits-and-generics/outputs/shapes_dyn.txt}{Each \code{area} call reaches the right
implementation.}\label{out:b2:rust-traits-and-generics:shapes}

\code{fat_pointer.rs} prints the sizes of a few pointers and reads the two halves of three trait-object references.

\omcode{code/bachelor-2/08-rust-traits-and-generics/listings/fat_pointer.rs}{12}{16}{The two words of a
\code{&dyn Shape}}\label{lst:b2:rust-traits-and-generics:fat}

\omoutput{code/bachelor-2/08-rust-traits-and-generics/outputs/fat_pointer.txt}{Sizes in bytes, on a target with
8-byte pointers.}\label{out:b2:rust-traits-and-generics:fat}

\begin{proposition}[A trait object reference is two pointers]\label{prop:b2:rust-traits-and-generics:fat}
A reference or box to a \code{dyn Trait} occupies two machine words, the data pointer and the \omterm{def:b2:cpp-polymorphism:vtable}{vtable} pointer. All
values of one type share one \omterm{def:b2:cpp-polymorphism:vtable}{vtable} for a given \omterm{def:b2:rust-traits-and-generics:trait}{trait}, and values of different types have different \omterm{def:b2:cpp-polymorphism:vtable}{vtables}.
\end{proposition}

\begin{proof}
A call \code{s.area()} on \code{s: &dyn Shape} must find, at run time, the code of \code{area} for the value's type,
and, to drop a \code{Box<dyn Shape>}, its drop function and size. These depend on the type, not on the value, so one
table per (type, \omterm{def:b2:rust-traits-and-generics:trait}{trait}) pair suffices and the compiler builds exactly one, a constant in the program; the value's
address is a second word, since the value itself may be anywhere. The value is not changed, so the table address
cannot be inside it: it travels in the pointer. Two words, as printed for \code{&dyn Shape} and \code{Box<dyn
Shape>}; the slices \code{&[u32]} and \code{&str} are \omterm{def:b2:rust-traits-and-generics:trait-object}{fat pointers} too, with a length as second word.
\end{proof}

\begin{omfigure}
\begin{tikzpicture}[font=\scriptsize, c/.style={draw, minimum width=17mm, minimum height=5mm}, v/.style={draw, fill=omDef!10, minimum width=26mm, minimum height=4.5mm}]
\node[font=\scriptsize\bfseries] at (1.6,1.0) {Rust: \texttt{\&dyn Shape}};
\node[c, fill=black!5] (d) at (0.75,0) {data};
\node[c, fill=black!5] (t) at (2.45,0) {vtable};
\node[c, fill=omThm!10, minimum width=22mm] (val) at (0.2,-1.6) {\texttt{Circle \{ r: 2.0 \}}};
\node[v] (v1) at (3.8,-1.3) {drop function};
\node[v, below=0pt of v1] (v2) {size 8, align 8};
\node[v, below=0pt of v2] (v3) {\texttt{area} of \texttt{Circle}};
\draw[ptr] (d.south) -- (val.north);
\draw[ptr] (t.south) -- (v1.north);
\node[text=black!60, align=center] at (3.8,-3.35) {one table per (type, trait),\\shared by all circles};
\begin{scope}[xshift=8.2cm]
\node[font=\scriptsize\bfseries] at (1.0,1.0) {C++: \texttt{Shape *}};
\node[c, fill=black!5] (p) at (0,0) {pointer};
\node[c, fill=black!5] (vp) at (0,-1.3) {vptr};
\node[c, fill=omThm!10, below=0pt of vp] (r) {\texttt{r = 2.0}};
\node[v] (w1) at (2.9,-1.3) {destructor};
\node[v, below=0pt of w1] (w2) {\texttt{area} of \texttt{Circle}};
\draw[ptr] (p.south) -- (vp.north);
\draw[ptr] (vp.east) -- (w1.west);
\node[text=black!60, align=center] at (1.2,-3.35) {the table pointer is inside\\every object};
\end{scope}
\end{tikzpicture}
\par\smallskip{\small Where the \omterm{def:b2:cpp-polymorphism:vtable}{table pointer} lives. In Rust it travels in the reference, so a \code{Circle}
stays 8 bytes and gains no hidden field; only the references used as \omterm{def:b2:rust-traits-and-generics:trait-object}{trait objects} grow, to 16 bytes. In C++ every
object of a class with \omterm{def:b2:cpp-polymorphism:virtual}{virtual functions} carries it, whether or not it is ever called virtually.}
\end{omfigure}

Not every \omterm{def:b2:rust-traits-and-generics:trait}{trait} can be used this way.

\begin{definition}[Object safety]\label{def:b2:rust-traits-and-generics:object-safe}
A \omterm{def:b2:rust-traits-and-generics:trait}{trait} is \emph{object safe}\index{object safety}, or \emph{dyn compatible}, if \code{dyn Trait} can be formed: its
methods can be called through a \omterm{def:b2:cpp-polymorphism:vtable}{vtable}. Methods that return \code{Self} by value, or have type parameters of their
own, break this (the caller does not know the size of \code{Self}; a generic method would need one table entry per
type argument), unless they are excluded with \code{where Self: Sized}.
\end{definition}

\omcode{code/bachelor-2/08-rust-traits-and-generics/listings/object_safety.rs}{3}{14}{A method that returns
\code{Self}}\label{lst:b2:rust-traits-and-generics:object-safety}

\omoutput[2-4,13-14]{code/bachelor-2/08-rust-traits-and-generics/outputs/object_safety.txt}{The message uses the
newer name, \emph{dyn compatible}.}\label{out:b2:rust-traits-and-generics:object-safety}

\begin{omfigure}
\begin{tikzpicture}[font=\scriptsize, b/.style={draw, rounded corners=2pt, fill=black!4, minimum width=32mm, minimum height=7mm, align=center}]
\node[b] (g) at (0,0) {\texttt{fn largest<T: PartialOrd>}\\ source, checked once};
\node[b, fill=omDef!10] (g1) at (5,0.7) {\texttt{largest::<u32>}\\ integer compare};
\node[b, fill=omDef!10] (g2) at (5,-0.7) {\texttt{largest::<f64>}\\ float compare};
\draw[ptr] (g.east) -- (g1.west); \draw[ptr] (g.east) -- (g2.west);
\node[text=black!60] at (5,1.6) {machine code};
\begin{scope}[yshift=-2.9cm]
\node[b] (h) at (0,0) {\texttt{fn sum\_dyn(\&[Box<dyn Shape>])}\\ source, checked once};
\node[b, fill=omDef!10] (h1) at (5,0) {\texttt{sum\_dyn}\\ one copy, calls via table};
\node[b, fill=omThm!10, minimum width=26mm] (t1) at (9.7,0.45) {\texttt{Rect}'s \texttt{area}};
\node[b, fill=omThm!10, minimum width=26mm] (t2) at (9.7,-0.45) {\texttt{Square}'s \texttt{area}};
\draw[ptr] (h.east) -- (h1.west);
\draw[ptr, dashed] (h1.east) -- (t1.west); \draw[ptr, dashed] (h1.east) -- (t2.west);
\node[text=black!60] at (9.7,1.15) {chosen at run time};
\end{scope}
\end{tikzpicture}
\par\smallskip{\small Static and \omterm{def:b2:cpp-polymorphism:dispatch}{dynamic dispatch}. Above, a generic is copied per type and each copy calls its
methods directly, where they can be inlined. Below, a function on \omterm{def:b2:rust-traits-and-generics:trait-object}{trait objects} exists once, and each call goes
through the \omterm{def:b2:cpp-polymorphism:vtable}{vtable} of the value at hand.}
\end{omfigure}

The two costs can be measured. \code{dispatch_cost.rs} sums the areas of a million shapes, twenty times, through a
\omterm{def:b2:rust-traits-and-generics:bound}{generic function} over rectangles only, through a \code{match} on an enum of the two shapes
(\cref{ch:b2:rust-enums-and-errors}), and through boxed \omterm{def:b2:rust-traits-and-generics:trait-object}{trait objects}.

\omcode{code/bachelor-2/08-rust-traits-and-generics/listings/dispatch_cost.rs}{13}{19}{The three
sums}\label{lst:b2:rust-traits-and-generics:dispatch}

\omoutput{code/bachelor-2/08-rust-traits-and-generics/outputs/run_dispatch.txt}{Measured on the author's laptop, two
runs; the times vary from run to run.}\label{out:b2:rust-traits-and-generics:dispatch}

The generic loop is the fastest: no call at all remains, and the additions are vectorised. The enum costs a test
per shape; \omterm{def:b2:rust-traits-and-generics:trait-object}{trait objects} cost an indirect call, and a pointer to follow to each box, about twice the enum. A few
nanoseconds per call matters only in inner loops: \omterm{def:b2:rust-traits-and-generics:trait-object}{trait objects} are the right tool for a list of plugins, handlers
or shapes of open-ended kinds, an enum for a closed set of kinds, a generic for code that must be fast for each
type.

\section{The standard traits}\label{sec:b2:rust-traits-and-generics:std}

A handful of \omterm{def:b2:rust-traits-and-generics:trait}{traits} of the standard library give a type its everyday abilities; most can be derived.

\begin{center}\footnotesize
\begin{tabular}{@{}lll@{}}
\hline
\omterm{def:b2:rust-traits-and-generics:trait}{trait} & gives & derive? \\
\hline
\code{Debug}, \code{Display} & \code{{:?}} and \code{{}} printing & \code{Debug} only \\
\code{Clone}, \code{Copy} & \code{.clone()}; copy on assignment instead of move & yes \\
\code{PartialEq}, \code{Eq} & \code{==}; \code{Eq} promises reflexivity, $x = x$ & yes \\
\code{PartialOrd}, \code{Ord} & \code{<}, \code{partial_cmp}; \code{Ord}: a total order, \code{cmp}, \code{sort} & yes \\
\code{Hash} & a key of \code{HashMap} and \code{HashSet} & yes \\
\code{Default} & \code{T::default()} & yes \\
\code{From}, \code{Into} & conversions, used by \code{?} & no \\
\code{Add}, \code{Mul}, \code{Neg}, \dots & the operators \code{+}, \code{*}, unary \code{-} & no \\
\code{Iterator} & \code{for} loops, \code{map}, \code{sum} (\cref{ch:b2:functional-programming}) & no \\
\code{Drop} & code run when the value is dropped & no \\
\hline
\end{tabular}
\end{center}

\begin{omfigure}
\begin{tikzpicture}[font=\scriptsize, t/.style={draw, rounded corners=2pt, fill=omDef!10, minimum width=20mm, minimum height=6mm}]
\node[t] (pe) at (0,0) {\texttt{PartialEq}};
\node[t] (eq) at (-2.6,-1.3) {\texttt{Eq}};
\node[t] (po) at (2.6,-1.3) {\texttt{PartialOrd}};
\node[t] (o) at (0,-2.6) {\texttt{Ord}};
\draw[ptr] (pe) -- (eq); \draw[ptr] (pe) -- (po); \draw[ptr] (eq) -- (o); \draw[ptr] (po) -- (o);
\node[anchor=west, text=black!70, align=left] at (4.3,-0.4) {\texttt{f64}: \texttt{PartialEq}, \texttt{PartialOrd} only\\(\texttt{NaN != NaN}, and \texttt{NaN} is unordered)};
\node[anchor=west, text=black!70, align=left] at (4.3,-2.2) {\texttt{u32}, \texttt{String}, \texttt{char}, tuples and\\derived structs of these: all four};
\end{tikzpicture}
\par\smallskip{\small The comparison \omterm{def:b2:rust-traits-and-generics:trait}{traits}. An arrow goes from a \omterm{def:b2:rust-traits-and-generics:trait}{trait} to one that requires it: a type can
implement \code{Ord} only if it implements \code{Eq} and \code{PartialOrd}.}
\end{omfigure}

Derived comparisons follow the order in which the fields are written.

\begin{proposition}[Derived orders are lexicographic]\label{prop:b2:rust-traits-and-generics:derived-ord}
For a struct with fields $f_1, \dots, f_n$ in this order, the derived \code{Ord} compares $a$ and $b$ by their first
fields, and if equal by their second, and so on: the lexicographic order on $(a.f_1, \dots, a.f_n)$. For an enum, a
value of an earlier variant is smaller than one of a later variant, and two values of the same variant compare
lexicographically on their fields. If the order of every field type is total, the derived order is total.
\end{proposition}

\begin{proof}
The first two claims describe the derived code, which is \code{a.f1.cmp(&b.f1)} followed, while the result is
\code{Equal}, by the next fields; for an enum, the comparison of the variants' indices comes first. Totality, for a
struct, by induction on $n$. For $n = 0$ all values are equal. For $n \ge 1$: if $a.f_1 \ne b.f_1$, the total order
on $f_1$ puts one first, and so does the derived order; if $a.f_1 = b.f_1$, the derived order is that of the
remaining $n - 1$ fields, total by the induction hypothesis. Transitivity and antisymmetry follow in the same way
(\cref{exo:b2:rust-traits-and-generics:8}). For an enum, the variant indices are integers, totally ordered, and within
one variant the struct case applies.
\end{proof}

\omcode{code/bachelor-2/08-rust-traits-and-generics/listings/people_sort.rs}{4}{24}{Derived orders on a person and an
age group}\label{lst:b2:rust-traits-and-generics:people}

\omoutput[2-7]{code/bachelor-2/08-rust-traits-and-generics/outputs/run_people.txt}{The town's 120 people, by name, by
age, and counted by group.}\label{out:b2:rust-traits-and-generics:people}

\code{f64} has no \code{Ord}: \code{NaN} compares neither smaller, equal nor greater than anything, itself included,
so a vector of \code{f64} cannot call \code{sort}.

\omoutput{code/bachelor-2/08-rust-traits-and-generics/outputs/float_sort.txt}{\code{float_sort.rs}: the comparisons
of NaN, then three ways to sort \code{[4.3, -0.1, NaN, 12.5, 3.9, -2.0]}.}\label{out:b2:rust-traits-and-generics:float}

\code{f64::total_cmp} is a total order on all bit patterns, \code{NaN} after $+\infty$. Treating \code{NaN} as
equal to everything gives a relation that is not transitive ($-0.1 = \mathrm{NaN} = 12.5$, yet $-0.1 < 12.5$), and
the sort, whose correctness assumes an order, returns a vector that is not sorted. The \code{unwrap} at least stops.

Operators are \omterm{def:b2:rust-traits-and-generics:trait}{traits} too. \code{a + b} is \code{Add::add(a, b)}, so implementing \code{Add} for a type gives it
\code{+}.

\begin{definition}[Associated type]\label{def:b2:rust-traits-and-generics:associated}
An \emph{associated type}\index{associated type} is a type named inside a \omterm{def:b2:rust-traits-and-generics:trait}{trait}, \code{type Output;}, that each
implementation fixes. \code{Add} declares \code{type Output}, the type of \code{a + b}; \code{Iterator} declares
\code{type Item}, the type of the elements it yields.
\end{definition}

\omcode{code/bachelor-2/08-rust-traits-and-generics/listings/vector2.rs}{10}{23}{Plane vectors with \code{+},
\code{-} and \code{*}}\label{lst:b2:rust-traits-and-generics:vector}

\omoutput{code/bachelor-2/08-rust-traits-and-generics/outputs/vector2.txt}{Walking halfway from the library to the
station.}\label{out:b2:rust-traits-and-generics:vector}

An \omterm{def:b2:rust-traits-and-generics:associated}{associated type} differs from a type parameter in who chooses it: \code{Mul<f64> for V2} and a possible
\code{Mul<V2> for V2} are two implementations, selected by the right operand's type, but each fixes its own
\code{Output}, so the type of \code{a * b} is never ambiguous.

\section{Newtypes and zero cost}\label{sec:b2:rust-traits-and-generics:zero}

\begin{definition}[Newtype]\label{def:b2:rust-traits-and-generics:newtype}
A \emph{newtype}\index{newtype} is a struct with a single field, \code{struct Metres(f64);}, that gives an existing
type a new name and a new set of operations: those implemented for it, and no others.
\end{definition}

A \omterm{def:b2:rust-traits-and-generics:newtype}{newtype} is used to give distinct meanings distinct types, a length and a time both stored as \code{f64}, and to
implement a foreign \omterm{def:b2:rust-traits-and-generics:trait}{trait} for a foreign type, which the orphan rule forbids directly
(\cref{exo:b2:rust-traits-and-generics:10} does it for \code{bool}).

\begin{definition}[Zero-cost abstraction]\label{def:b2:rust-traits-and-generics:zero-cost}
An abstraction of a language is \emph{zero-cost}\index{zero-cost abstraction} if a program written with it compiles
to code no slower and no larger than the program a careful programmer would write without it, by hand.
\end{definition}

Generics are \omterm{def:b2:rust-traits-and-generics:zero-cost}{zero-cost} in this sense (a copy per type, as if written by hand); \omterm{def:b2:rust-traits-and-generics:trait-object}{trait objects} are not quite, by
design: they buy the ability to mix types for an indirect call. \omterm{def:b2:rust-traits-and-generics:newtype}{Newtypes} are the clearest case, and the weekend
problem measures it.

\begin{inpractice}[Choosing between generics, enums and trait objects]
Take a generic parameter when the type is known at compile time at every call site; an enum when the set of kinds is
closed and you own it, so that a new kind is a compile error at every \code{match}; a \omterm{def:b2:rust-traits-and-generics:trait-object}{trait object} when the set is
open, other crates add kinds, or the values must be mixed in one collection. Derive \code{Debug} on every type, and
\code{Clone}, \code{PartialEq}, \code{Eq}, \code{Hash}, \code{Ord} where they make sense.
\end{inpractice}

\section{Exercises}\label{sec:b2:rust-traits-and-generics:exercises}

\begin{exercise}[$\star$]\label{exo:b2:rust-traits-and-generics:1}
Predict the five lines. Which \code{name} does each call reach, and why do the second and fourth lines differ from
what one might expect?
\omcode{code/bachelor-2/08-rust-traits-and-generics/listings/exo_which.rs}{12}{33}{\omterm{def:b2:rust-traits-and-generics:trait}{Trait} methods, inherent methods,
generics and \code{dyn}}
\end{exercise}

\begin{exercise}[$\star$]\label{exo:b2:rust-traits-and-generics:2}
\code{report} does not compile, with two errors. Find the bounds it needs, and write its signature with a
\code{where} clause.
\omcode{code/bachelor-2/08-rust-traits-and-generics/listings/exo_bound.rs}{5}{15}{Printing and counting}
\end{exercise}

\begin{exercise}[$\star$]\label{exo:b2:rust-traits-and-generics:3}
Write a type \code{Fraction} that is always in lowest terms with a positive denominator, a \omterm{def:b2:cpp-classes-raii:constructor}{constructor} returning
\code{None} for a zero denominator, and \code{Display} printing \code{3/4}, \code{-3/4}, or \code{2} for a whole
number. Make \code{format!("{:>5}", f)} pad the fraction. Why does the derived \code{PartialEq} then compare values
correctly?
\end{exercise}

\begin{exercise}[$\star$]\label{exo:b2:rust-traits-and-generics:4}
A program has \code{fn show<T: Display>(x: T)} and \code{fn pair<A: Display, B: Display>(a: A, b: B)}, whose body
calls \code{show(a)} and \code{show(b)}. \code{main} calls \code{show(3)}, \code{show("Mill")}, \code{show(4)},
\code{pair(1, "x")} and \code{pair("y", 2.5)}. How many copies of \code{show} and of \code{pair} does the compiled
program contain (before inlining)? Justify with \cref{prop:b2:rust-traits-and-generics:mono}.
\end{exercise}

\begin{exercise}[$\star\star$]\label{exo:b2:rust-traits-and-generics:5}
Predict the two sorted lists, using \cref{prop:b2:rust-traits-and-generics:derived-ord}. Then reorder the fields of
\code{Loan} so that loans sort by title first: what else changes?
\omcode{code/bachelor-2/08-rust-traits-and-generics/listings/exo_derived.rs}{2}{27}{An enum and a struct, sorted}
\end{exercise}

\begin{exercise}[$\star\star$]\label{exo:b2:rust-traits-and-generics:6}
The median of seven temperatures comes out as $-0.4$. Explain, using the output, why the sort did not sort, and fix
the program twice: with \code{total_cmp}, and by refusing \code{NaN} when the fields are parsed.
\omcode{code/bachelor-2/08-rust-traits-and-generics/listings/exo_nan.rs}{4}{14}{A median}
\omoutput[2-3]{code/bachelor-2/08-rust-traits-and-generics/outputs/exo_nan.txt}{The result and the vector.}
\end{exercise}

\begin{exercise}[$\star\star$]\label{exo:b2:rust-traits-and-generics:7}
Write \code{largest<T: PartialOrd + Copy>(xs: &[T]) -> Option<T>}, then \code{largest_ref}, which needs no
\code{Copy}. What does each return on \code{[1.0, NaN, 3.0]} and on \code{[NaN, 3.0]}, and why?
\end{exercise}

\begin{exercise}[$\star\star$]\label{exo:b2:rust-traits-and-generics:8}
Complete the proof of \cref{prop:b2:rust-traits-and-generics:derived-ord}: prove that the derived order of a struct
whose fields have total orders is transitive and antisymmetric.
\end{exercise}

\begin{exercise}[$\star\star$]\label{exo:b2:rust-traits-and-generics:9}
Write a \omterm{def:b2:rust-traits-and-generics:trait}{trait} \code{Shape} with \code{area}, \code{perimeter} and \code{name}, three shapes, and a function that
sorts a \code{Vec<Box<dyn Shape>>} by decreasing area and returns the names. Why can you not derive \code{Ord} on
\code{Box<dyn Shape>}, and what do you use instead?
\end{exercise}

\begin{exercise}[$\star\star\star$]\label{exo:b2:rust-traits-and-generics:10}
Write a \omterm{def:b2:rust-traits-and-generics:trait}{trait} \code{Numeric} (\code{Copy}, \code{+}, \code{*}, a zero and a one) and a square matrix type
generic over it, with \code{identity}, \code{mul} and \code{pow} by repeated squaring. Implement \code{Numeric} for
\code{i64}, \code{f64}, and for booleans with \emph{or} as addition and \emph{and} as multiplication: why must you
use a \omterm{def:b2:rust-traits-and-generics:newtype}{newtype} for the last one? Check that the tenth power of $\bigl(\begin{smallmatrix}1&1\\1&0\end{smallmatrix}
\bigr)$ contains the Fibonacci number 55, and what the boolean powers of a graph's adjacency matrix tell.
\end{exercise}

\begin{exercise}[$\star\star\star$]\label{exo:b2:rust-traits-and-generics:11}
A program has $m$ types, each implementing $k$ \omterm{def:b2:rust-traits-and-generics:trait}{traits}, the \omterm{def:b2:rust-traits-and-generics:trait}{trait} $i$ having $n_i$ methods, all used as \omterm{def:b2:rust-traits-and-generics:trait-object}{trait
objects}. (a) How many \omterm{def:b2:cpp-polymorphism:vtable}{vtables} does it hold, and how many words in all, if each table has three words (drop, size,
alignment) plus one per method? (b) A C++ program with the same $m$ classes, each deriving from the $k$ abstract
bases: how many \omterm{def:b2:cpp-polymorphism:vtable}{table pointers} does each object carry? (c) Which design costs memory per object, and which per type?
\end{exercise}

\begin{exercise}[$\star\star\star$]\label{exo:b2:rust-traits-and-generics:12}
Write \code{Connection<State>}, where \code{State} is \code{Closed} or \code{Open}, two empty types: \code{new}
gives a closed connection, \code{open} consumes it and returns an open one, \code{send} exists only on
\code{Connection<Open>}, and \code{close} returns a closed one. Show that sending on a closed connection, or using a
connection after closing it, does not compile, and that the state costs no byte.
\end{exercise}

\section{Problem: A units library}\label{sec:b2:rust-traits-and-generics:problem}

\begin{problem}[{Weekend problem --- A units library}]\label{pb:b2:rust-traits-and-generics:1}
% ledger: b2:mco-units
% ledger: b2:mco-loss
On 23 September 1999 a spacecraft sent to orbit Mars was lost as it arrived. The investigation board found that a
ground program wrote the impulses of the thrusters in pound-force seconds, where the interface specification required
newton-seconds; the navigation software read the numbers as newton-seconds, and underestimated the effect of the
thrusters by a factor of 4.45. The reconstructed trajectory passed 57\,km above the surface, too low to survive. Both
programs passed numbers, and a number does not carry its unit. This problem makes the unit part of the type and
checks what it costs.

\medskip\noindent\textbf{Part I --- Newtypes.}\par\nopagebreak
\omcode{code/bachelor-2/08-rust-traits-and-generics/listings/units.rs}{5}{27}{Three units and four
operations}
\begin{enumerate}[beginpenalty=10000]
\item What is the size of \code{Metres}? Of \code{(Metres, Seconds)}?
\item Why is the field \code{pub}, and what would be lost if it were private? What would be gained?
\item Which \omterm{def:b2:rust-traits-and-generics:trait}{traits} do the three types derive, and why not \code{Eq} or \code{Ord}?
\item Could the board's bug have been written with these types, if the impulse were a type \code{NewtonSeconds}?
What would the ground program have had to write?
\item Write a type \code{PoundForceSeconds} and a conversion \code{From<PoundForceSeconds> for NewtonSeconds}.
\item Why is a conversion function safer than a constant 4.45 multiplied in by each caller?
\end{enumerate}

\medskip\noindent\textbf{Part II --- Operator traits.}\par\nopagebreak
\omoutput{code/bachelor-2/08-rust-traits-and-generics/outputs/run_units.txt}{The library at work, and two mistakes
refused (\code{units_mix.rs}).}
\begin{omfigure}
\begin{tikzpicture}[font=\scriptsize, u/.style={draw, rounded corners=2pt, fill=omDef!10, minimum width=26mm, minimum height=6mm}]
\node[u] (m) at (0,0) {\texttt{Metres}};
\node[u] (s) at (6,0) {\texttt{Seconds}};
\node[u] (v) at (3,-1.8) {\texttt{MetresPerSecond}};
\draw[ptr] (m.south) -- node[left, xshift=-2pt] {\texttt{/ Seconds}} (v.north west);
\draw[ptr] (v.north east) -- node[right, xshift=2pt] {\texttt{* Seconds}: \texttt{Metres}} (s.south);
\draw[ptr] ([xshift=-5mm]m.north) .. controls +(-0.2,0.9) and +(0.2,0.9) .. node[above] {\texttt{+ Metres}} ([xshift=5mm]m.north);
\draw[ptr] ([xshift=-5mm]s.north) .. controls +(-0.2,0.9) and +(0.2,0.9) .. node[above] {\texttt{+ Seconds}} ([xshift=5mm]s.north);
\node[omThm, align=left, anchor=west] at (8.0,-0.6) {refused:\\ \texttt{Metres + Seconds}\\ \texttt{let x: f64 = d;}\\ \texttt{Metres * Metres}};
\end{tikzpicture}
\par\smallskip{\small The operations the types allow. The arrow from \code{MetresPerSecond} leads, through a
multiplication by \code{Seconds}, back to \code{Metres}; every operation not drawn is a compile error.}
\end{omfigure}
\begin{enumerate}[resume, beginpenalty=10000]
\item Which impl makes \code{d / t} compile, and what is the type of the result?
\item Why does \code{Metres * Metres} not compile? Should it, and with which output type?
\item The two errors of \code{units_mix.rs} are both \emph{mismatched types}. Explain each.
\item \code{v * Seconds(3600.0)} works but \code{Seconds(3600.0) * v} does not. Add the missing impl.
\item Implement \code{std::iter::Sum} for \code{Metres}, so that \code{legs.iter().map(|l| l.0).sum()} gives
\code{Metres}.
\item Why is \code{Display} implemented for \code{MetresPerSecond} but \code{Debug} printed for the others in the
output?
\end{enumerate}

\medskip\noindent\textbf{Part III --- Generics over units.}\par\nopagebreak
\begin{enumerate}[resume, beginpenalty=10000]
\item Write \code{fn total<U: Add<Output = U> + Copy>(xs: &[U], zero: U) -> U}, and use it on lengths and on times.
\item How many copies of \code{total} does your program then contain?
\item Can \code{total} be called on a slice that mixes \code{Metres} and \code{Seconds}? Why is that a feature?
\item A \omterm{def:b2:rust-traits-and-generics:trait}{trait} \code{Unit} with an associated constant \code{SYMBOL: &'static str} lets one generic \code{Display}
serve every unit. Sketch it.
\item What would a \omterm{def:b2:rust-traits-and-generics:trait-object}{trait object} \code{&dyn Unit} cost compared with the generic?
\item Name two units that a real library would add, and the operations between them.
\end{enumerate}

\medskip\noindent\textbf{Part IV --- Zero cost.}\par\nopagebreak
\omcode{code/bachelor-2/08-rust-traits-and-generics/listings/units_asm.rs}{6}{24}{The same trip, typed and
untyped}
\omoutput{code/bachelor-2/08-rust-traits-and-generics/outputs/run_asm.txt}{The object file's symbols, and the size
of the code.}
\begin{enumerate}[resume, beginpenalty=10000]
\item What does the attribute \code{no_mangle} change, and why is it used here?
\item Read the two symbol lines: what do the section, the size and the address say about the two functions?
\item How many instructions does the function have, and what are its last three doing?
\item Nine \code{addpd} instructions add two \code{f64} at once. What is being added in pairs?
\item What would have to differ in the typed version for the two to compile differently?
\item State the result of the problem.
\end{enumerate}
\end{problem}
